\documentclass[12pt]{article}
\usepackage[paperwidth=58.8cm,voffset=-11pt,paperheight=33cm,margin=1cm]{geometry}

%======================================================================
% PAQUETES
%======================================================================
\usepackage[spanish]{babel}
\usepackage{anyfontsize}
\usepackage{enumitem}
\usepackage{ifthen}
\usepackage{multido}
\usepackage[]{textpos}
\usepackage{xcolor}
\usepackage[most]{tcolorbox}
\usetikzlibrary{calc,patterns,decorations.pathmorphing}
\usepackage[explicit]{titlesec}
\usepackage{graphicx}

%======================================================================
% IMAGENES
% Los diagramas viven en imgs/ como .svg. LuaLaTeX no lee SVG, asi que
% latexmk (ver .latexmkrc) los convierte a PDF vectorial en build/svg/
% manteniendo la misma estructura de carpetas. Por eso las secciones
% escriben la ruta del .svg sin extension y graphicx la resuelve aqui.
%======================================================================
\graphicspath{{build/svg/}}

%======================================================================
% TIPOGRAFIA
% Montserrat para titulos (geometrica, pesa bien en formato poster),
% Lato para el texto corrido (humanista, comoda de leer en parrafo).
% Ambas vienen con TeX Live: no hay que instalar nada en el sistema.
%======================================================================
\usepackage{fontspec}
\setmainfont{Lato}[
  Extension      = .ttf,
  UprightFont    = *-Regular,
  BoldFont       = *-Bold,
  ItalicFont     = *-Italic,
  BoldItalicFont = *-BoldItalic,
]
\newfontfamily\FuenteTitulos{Montserrat}[
  Extension      = .otf,
  UprightFont    = *-SemiBold,
  BoldFont       = *-Bold,
  ItalicFont     = *-SemiBoldItalic,
  BoldItalicFont = *-BoldItalic,
]

\usepackage[backend=biber,sorting=none]{biblatex}
\addbibresource{referencias/Textos.bib}
\addbibresource{referencias/Articulos.bib}
\addbibresource{referencias/Libros.bib}

%======================================================================
% PALETA
%======================================================================
\colorlet{CLinea}{gray!70!white}
\colorlet{CFondo}{gray!5!white}
\colorlet{CBlanco}{white}
\colorlet{CTSeccion}{blue!30}
\definecolor{CTDiario}{RGB}{133, 193, 233}
\definecolor{CCDiario}{RGB}{253, 242, 233}
\definecolor{CCPoster}{RGB}{253, 254, 254}

% Trazo "a mano alzada" de los marcos
\tikzset{decoration={random steps,segment length=5mm,amplitude=0.3pt}}

%======================================================================
% DATOS DEL POSTER
%======================================================================
\newcommand{\MATERIA}{ARQUITECTURA OmniStream}
\newcommand{\Ciudad}{Bogotá D.C.}
\newcommand{\Escudo}{imgs/stjCartoon}

\setlength\parindent{0pt}

%======================================================================
% GRILLA DE MAQUETACION
% La pagina se divide en filas x columnas; textpos posiciona cada bloque
% en coordenadas de esa grilla.
%======================================================================
\NewDocumentCommand{\DibujarLinea}{m m m m m}{
  \setlength{\parindent}{0cm}
  \setlength{\parskip}{0cm}
  \count3=#1
  \advance\count3 by 1
  \multido{\i=0+1}{\count3}{
    \ifthenelse{\equal{#5}{horizontal}}{
      \begin{textblock}{#2}(0,\i)
        \rule[0pt]{#3}{#4}
      \end{textblock}
    }{
      \begin{textblock}{#2}(\i,0)
        \rule[0pt]{#3}{#4}
      \end{textblock}
    }
  }
}

\newdimen\CeldaW
\newdimen\CeldaH
\newcount\nFilas
\newcount\nColus

\NewDocumentCommand{\CalcularGrilla}{m m O{\textwidth} O{\textheight} O{show}}{
  \nFilas=#1
  \nColus=#2

  \CeldaW=#3
  \CeldaH=#4
  \divide\CeldaH by \nFilas
  \divide\CeldaW by \nColus

  \setlength{\TPHorizModule}{\CeldaW}
  \setlength{\TPVertModule}{\CeldaH}

  \ifthenelse{\equal{#5}{show}}{
    \DibujarLinea{\nFilas}{\nColus}{\textwidth}{.1pt}{horizontal}
    \DibujarLinea{\nColus}{\nFilas}{.1pt}{\textheight}{vertical}
  }{}
}

\NewDocumentEnvironment{NuevaPagina}{m m m}
{
  \thispagestyle{empty}
  \clearpage
  \CalcularGrilla{#1}{#2}[\textwidth][\textheight][#3]
}
{
  \newpage
}

\newdimen\Margen
\setlength{\Margen}{2pt}
\newdimen\AltoP
\NewDocumentEnvironment{NuevoParrafo}{m m O{1} O{1}}
{
  \setlength{\AltoP}{\dimexpr#4\TPVertModule-2\Margen\relax}
  \TPMargin{\Margen}
  \begin{textblock}{#3}(#2,#1)
}
{
  \end{textblock}
}

%======================================================================
% MARCOS
% Los cuatro primeros argumentos son el grosor (en pt) de cada borde;
% 0 = sin borde. El quinto es el color de fondo.
%======================================================================
\newcommand{\LineaIzqC}{0}
\newcommand{\LineaIzqU}{1}

\newcommand{\LineaDerC}{0}
\newcommand{\LineaDerU}{1}
\newcommand{\LineaDerD}{2}

\newcommand{\LineaSupC}{0}
\newcommand{\LineaSupU}{1}

\newcommand{\LineaInfC}{0}
\newcommand{\LineaInfU}{1}
\newcommand{\LineaInfD}{2}

\tcbset{
  codigoMarco/.style n args={5}{
    empty,
    height=\AltoP,
    frame code={\path[fill=#5](frame.south west) rectangle (frame.north east);};
    \ifthenelse{\equal{#1}{0}}{}{\draw[decorate,color=CLinea,line width =#1 pt](frame.north east) -- (frame.north west);};
    \ifthenelse{\equal{#2}{0}}{}{\draw[decorate,color=CLinea,line width =#2 pt](frame.south east) -- (frame.south west);};
    \ifthenelse{\equal{#3}{0}}{}{\draw[decorate,color=CLinea,line width =#3 pt](frame.north west) -- (frame.south west);};
    \ifthenelse{\equal{#4}{0}}{}{\draw[decorate,color=CLinea,line width =#4 pt](frame.north east) -- (frame.south east);};
  }
}
\tcbset{
  codigoLinea/.style n args={1}{
    empty,
    height=\AltoP,
    frame code={\path[fill=white](frame.south west) rectangle (frame.north east);};
    {\draw[decorate,color=#1,line width =10pt]([yshift=-12pt]frame.north east) -- ([yshift=-12pt]frame.north west);};
  }
}

\newtcolorbox{Linea}[1]{codigoLinea={#1}}

\newtcolorbox{MarcoSinTitulo}[5]{codigoMarco={#1}{#2}{#3}{#4}{#5}}
\NewDocumentEnvironment{Marco}{O{1} O{1} O{1} O{1} O{CFondo}}
{
  \begin{MarcoSinTitulo}{#1}{#2}{#3}{#4}{#5}
}
{
  \end{MarcoSinTitulo}
}

%======================================================================
% TITULOS
%======================================================================
\newcommand{\COLORSS}{\color{white}}
\titleformat{\subsection}{\centering\bf\LARGE\COLORSS}{\filcenter\LARGE\bf\FuenteTitulos\Alph{subsection}.#1}{0.0em}{}

\newcommand{\subseccionC}[1]{
  {\centering\FuenteTitulos\large\color{CTSeccion} \textbf{#1} \\ \vspace*{3pt} }
}

% Un diagrama dentro de un Marco. Se escala al maximo que quepa en la
% celda sin deformarse, descontando lo que ocupa el titulo. Asi no hay
% que calibrar un scale= a mano por figura ni se sale del recuadro si
% alguien reexporta el .svg con otro tamano.
\newcommand{\AltoTitulo}{34pt}
\newcommand{\Diagrama}[1]{%
  \centering\includegraphics[
    width=\linewidth,
    height=\dimexpr\AltoP-\AltoTitulo\relax,
    keepaspectratio]{#1}%
}

\newcommand{\CabeceraDeSubSeccionConAnchoN}[7]{
  \begin{NuevoParrafo}{#1}{#2}[16][2]
    \begin{Marco}[\LineaSupC][\LineaInfC][\LineaIzqC][\LineaDerC][CBlanco]
      \begin{center}
        \vspace{-5pt}
        \renewcommand{\COLORSS}{\color{#6}}
        \subsection{\textbf{#5}}
      \end{center}
    \end{Marco}
  \end{NuevoParrafo}
  \begin{NuevoParrafo}{#1}{\the\numexpr #2 + #3 \relax}[#4][1]
    \begin{Linea}{#7}
    \end{Linea}
  \end{NuevoParrafo}
}

%======================================================================
% CABECERA COMUN DEL POSTER
%======================================================================
\newcommand{\CabeceraDePoster}[1]{
  \begin{NuevoParrafo}{0}{0}[8][6]
    \begin{Marco}[\LineaSupC][\LineaInfC][\LineaIzqC][\LineaDerC][CCDiario]
      \includegraphics[width=65pt]{\Escudo}
    \end{Marco}
  \end{NuevoParrafo}
  \begin{NuevoParrafo}{0}{8}[72][6]
    \begin{Marco}[\LineaSupC][\LineaInfC][\LineaIzqC][\LineaDerC][CCDiario]
      \vspace{15pt}
      \centering\FuenteTitulos\fontsize{50}{50}\selectfont\color{CTDiario}\textbf{#1}
    \end{Marco}
  \end{NuevoParrafo}
}

\newcommand{\CiudadFechaVolumen}{
  \begin{NuevoParrafo}{6}{0}[94][2]
    \begin{Marco}[\LineaSupU][\LineaInfU][\LineaIzqC][\LineaDerC][CBlanco]
      \centering\textbf{--- \Ciudad\ \today\ ---}
    \end{Marco}
  \end{NuevoParrafo}
}

\newcommand{\autores}[3]{
  \vspace{-5pt}
  \large \centering{\FuenteTitulos STJ}
  \vspace{5pt}
  \hrule
  \vspace{5pt}
  OmniStream Team:\\ #1\\#2\\#3
}

\newenvironment{MisionVision}[1][Mision]{
  \begin{minipage}[t][.20\textheight]{\textwidth}
    \begin{center}
      \color{gray}\FuenteTitulos\bfseries{\fontsize{20}{10}\selectfont #1}
      \vspace{-10pt}
    \end{center}
}
{
  \end{minipage}\\[1cm]
}

\newenvironment{Objetivos}[1][Objetivos]{
  \begin{minipage}[t][.20\textheight]{\textwidth}
    \begin{center}
      \color{gray}\FuenteTitulos\bfseries{\fontsize{20}{10}\selectfont #1}
      \vspace{-10pt}
    \end{center}
}
{
  \end{minipage}\\[1cm]
}

\renewcommand{\MATERIA}{ARQUITECTURA OmniStream}
\renewcommand{\Ciudad}{Bogotá D.c}

\newcommand{\IntegrantesA}{***}
\newcommand{\IntegrantesB}{***}

\newcommand{\CabeceraComun}{
	\CabeceraDePoster{\MATERIA} 
	\begin{NuevoParrafo}{0}{80}[14][6]
		\begin{Marco}[\LineaSupC][\LineaInfC][\LineaIzqC][\LineaDerC][CBlanco]
			\autores{\IntegrantesA}{\IntegrantesB}
		\end{Marco}
	\end{NuevoParrafo} 
	\CiudadFechaVolumen
}

\newcommand{\Portada}{
	\CabeceraDeSubSeccionConAnchoN{9}{0}{14}{80}{OmniStream .CO}{CTDiario}{CCDiario}
	\begin{NuevoParrafo}{11}{0}[34][6]
		\begin{Marco}[\LineaSupC][\LineaInfC][\LineaIzqC][\LineaDerC][CBlanco]
			\begin{MisionVision}[Misión]
				{OmniStream proporciona a personas y organizaciones una plataforma auto-hospedada para almacenar, categorizar, consultar y reproducir sus recursos multimedia de forma autónoma, segura y principalmente offline, manteniendo el control total de su información.}
			\end{MisionVision}			
		\end{Marco}
	\end{NuevoParrafo}
	\begin{NuevoParrafo}{17}{0}[34][6]
		\begin{Marco}[\LineaSupC][\LineaInfC][\LineaIzqC][\LineaDerC][CBlanco]
			\begin{MisionVision}[Visión]
				{OmniStream se proyecta a 3 años como una solución de referencia en gestión multimedia auto-hospedada, reconocida por su modularidad, su experiencia de usuario multiplataforma y su independencia de servicios externos.}
			\end{MisionVision}
		\end{Marco}
	\end{NuevoParrafo}
		\begin{NuevoParrafo}{11}{38}[56][12]
		\begin{Marco}[\LineaSupC][\LineaInfC][\LineaIzqC][\LineaDerC][CBlanco]
			\begin{Objetivos}[Objetivos]
				{	{\large \textbf{General}:}\vspace{5pt}\\
					{\linespread{.8}\selectfont Diseñar e implementar un sistema modular de gestión multimedia que permita almacenar, organizar, consultar y reproducir contenido digital de manera eficiente y offline.}\\[0.3cm]
					{\large \textbf{Específicos:}}
					{\begin{itemize}
							\setlength{\itemsep}{-0.5pt}
							\setlength{\parskip}{0pt}
							\setlength{\parsep}{0pt}
							\item Implementar un módulo de almacenamiento e indexación que procese múltiples formatos de video, imagen, audio y texto extrayendo sus metadatos automáticamente.
							\item Desarrollar un sistema de categorización, búsqueda y filtrado adaptable a los criterios de cada usuario u organización.
							\item Exponer servicios de reproducción multiplataforma (Web, Desktop y Mobile) con control de acceso por usuario y funcionamiento sin conexión permanente a Internet. 
					\end{itemize}}
				}
			\end{Objetivos}
		\end{Marco}
	\end{NuevoParrafo}
}


%======================================================================
\begin{document}

\pgfmathsetseed{2024}

\begin{NuevaPagina}{64}{94}{show1}
	\definecolor{CCDiario}{HTML}{e3e8ff}
	\CabeceraComun
	\Portada
	\definecolor{CCPoster}{HTML}{7e7dd9}	
	\CabeceraDeSubSeccionConAnchoN{23}{0}{14}{80}{MOTIVACIÓN}{CTDiario}{CCPoster}
	\begin{NuevoParrafo}{25}{0}[40][16]
		\begin{Marco}[\LineaSupC][\LineaInfC][\LineaIzqC][\LineaDerC][CBlanco]
			\subseccionC{MetaModelo}
			\Diagrama{imgs/02-motivacion/Meta}
		\end{Marco}
	\end{NuevoParrafo}
	
	\begin{NuevoParrafo}{25}{41}[26][18]
		\begin{Marco}[\LineaSupC][\LineaInfC][\LineaIzqC][\LineaDerC][CBlanco]
		\subseccionC{Stakeholder}
		\Diagrama{imgs/02-motivacion/Interesado}
		\end{Marco}
	\end{NuevoParrafo}
	
	\begin{NuevoParrafo}{25}{68}[26][18]
		\begin{Marco}[\LineaSupC][\LineaInfC][\LineaIzqC][\LineaDerC][CBlanco]
			\subseccionC{Principios}
			\Diagrama{imgs/02-motivacion/Principios}
		\end{Marco}
	\end{NuevoParrafo}
	
	\begin{NuevoParrafo}{40}{3}[19][9]
		\begin{Marco}[\LineaSupC][\LineaInfC][\LineaIzqC][\LineaDerC][CBlanco]
			\subseccionC{Realización de Objetivos}
			\Diagrama{imgs/02-motivacion/RealObjetivos}
		\end{Marco}
	\end{NuevoParrafo}
		
	\begin{NuevoParrafo}{49}{3}[19][14]
		\begin{Marco}[\LineaSupC][\LineaInfC][\LineaIzqC][\LineaDerC][CBlanco]
			\subseccionC{Realización de Requerimientos}
			\Diagrama{imgs/02-motivacion/RealRequerimientos}
		\end{Marco}
	\end{NuevoParrafo}

	\begin{NuevoParrafo}{42}{27}[33][22]
		\begin{Marco}[\LineaSupC][\LineaInfC][\LineaIzqC][\LineaDerC][CBlanco]
			\subseccionC{Contribución de Objetivos}
			\Diagrama{imgs/02-motivacion/ContObjetivos}
		\end{Marco}
	\end{NuevoParrafo}
	\begin{NuevoParrafo}{42}{65}[26][20]
		\begin{Marco}[\LineaSupC][\LineaInfC][\LineaIzqC][\LineaDerC][CBlanco]
			\subseccionC{Motivación}
			\Diagrama{imgs/02-motivacion/Motivacion}
		\end{Marco}
	\end{NuevoParrafo}

\end{NuevaPagina}

\begin{NuevaPagina}{64}{94}{show1}
	\definecolor{CCDiario}{HTML}{fffec4}
	\CabeceraComun
	\definecolor{CCPoster}{HTML}{caaa9b}
	\CabeceraDeSubSeccionConAnchoN{9}{0}{14}{80}{ESTRATEGIA}{CTDiario}{CCPoster}
	\begin{NuevoParrafo}{11}{0}[23][12]
		\begin{Marco}[\LineaSupC][\LineaInfC][\LineaIzqC][\LineaDerC][CBlanco]
			\subseccionC{MetaModelo}
			\centering\includegraphics[scale=0.4]{\ImgEtg MetaModelo.pdf}	
		\end{Marco}
	\end{NuevoParrafo}	
	\begin{NuevoParrafo}{11}{25}[20][14]
		\begin{Marco}[\LineaSupC][\LineaInfC][\LineaIzqC][\LineaDerC][CBlanco]7
			\subseccionC{\PVMaC}
			\centering\includegraphics[scale=0.5]{\ImgEtg MapaCapacidad.pdf}		
		\end{Marco}
	\end{NuevoParrafo}	
	\begin{NuevoParrafo}{24}{25}[20][10]
		\begin{Marco}[\LineaSupC][\LineaInfC][\LineaIzqC][\LineaDerC][CBlanco]
			\subseccionC{\PVMaR}
			\centering\includegraphics[scale=0.5]{\ImgEtg MapaRecursos.pdf}		
		\end{Marco}
	\end{NuevoParrafo}	
		\begin{NuevoParrafo}{24}{45}[22][10]
		\begin{Marco}[\LineaSupC][\LineaInfC][\LineaIzqC][\LineaDerC][CBlanco]
			\subseccionC{\PVFdV}
			\centering\includegraphics[scale=0.5]{\ImgEtg FlujoValor.pdf}		
		\end{Marco}
	\end{NuevoParrafo}	
	\begin{NuevoParrafo}{11}{45}[22][13]
		\begin{Marco}[\LineaSupC][\LineaInfC][\LineaIzqC][\LineaDerC][CBlanco]
			\subseccionC{\PVRes}
			\centering\includegraphics[scale=0.45]{\ImgEtg RealizacionResultado.pdf}		
		\end{Marco}
	\end{NuevoParrafo}	

	\begin{NuevoParrafo}{11}{65}[29][23]
		\begin{Marco}[\LineaSupC][\LineaInfC][\LineaIzqC][\LineaDerC][CBlanco]
			\subseccionC{\PVEtg}
			\centering\includegraphics[scale=.5]{\ImgEtg Estrategia.pdf}	
		\end{Marco}
	\end{NuevoParrafo}

	\definecolor{CCPoster}{HTML}{f8f32b}	
	\CabeceraDeSubSeccionConAnchoN{34}{0}{14}{80}{NEGOCIO}{CTDiario}{CCPoster}
	\begin{NuevoParrafo}{36}{0}[40][15]
		\begin{Marco}[\LineaSupC][\LineaInfC][\LineaIzqC][\LineaDerC][CBlanco]
			\subseccionC{MetaModelo}
			\Diagrama{imgs/04-negocio/Meta}
		\end{Marco}
	\end{NuevoParrafo}
	
	\begin{NuevoParrafo}{36}{29}[20][15]
		\begin{Marco}[\LineaSupC][\LineaInfC][\LineaIzqC][\LineaDerC][CBlanco]
			\subseccionC{Organización}
			\Diagrama{imgs/04-negocio/Organizacion}
		\end{Marco}
	\end{NuevoParrafo}
	
	\begin{NuevoParrafo}{36}{45}[26][15]
		\begin{Marco}[\LineaSupC][\LineaInfC][\LineaIzqC][\LineaDerC][CBlanco]
			\subseccionC{Cooperación de Actor}
			\Diagrama{imgs/04-negocio/CoopActor}
		\end{Marco}
	\end{NuevoParrafo}
	
	\begin{NuevoParrafo}{36}{68}[26][15]
		\begin{Marco}[\LineaSupC][\LineaInfC][\LineaIzqC][\LineaDerC][CBlanco]
			\subseccionC{Función de Negocio}
			\Diagrama{imgs/04-negocio/FuncionNegocio}
		\end{Marco}
	\end{NuevoParrafo}
		
	\begin{NuevoParrafo}{51}{4}[30][13]
		\begin{Marco}[\LineaSupC][\LineaInfC][\LineaIzqC][\LineaDerC][CBlanco]
			\subseccionC{Proceso de Negocio}
			\Diagrama{imgs/04-negocio/ProcesoNegocio}
		\end{Marco}
	\end{NuevoParrafo}
	
	\begin{NuevoParrafo}{51}{36}[30][13]
		\begin{Marco}[\LineaSupC][\LineaInfC][\LineaIzqC][\LineaDerC][CBlanco]
			\subseccionC{Cooperación de Proceso de Negocio}
			\Diagrama{imgs/04-negocio/CoopProcesoNegocio}
		\end{Marco}
	\end{NuevoParrafo}
	
	\begin{NuevoParrafo}{51}{68}[26][13]
		\begin{Marco}[\LineaSupC][\LineaInfC][\LineaIzqC][\LineaDerC][CBlanco]
			\subseccionC{Producto}
			\Diagrama{imgs/04-negocio/Producto}
		\end{Marco}
	\end{NuevoParrafo}

\end{NuevaPagina}

\begin{NuevaPagina}{64}{94}{show1}
	\definecolor{CCDiario}{HTML}{e9fced}
	\CabeceraComun
	\definecolor{CCPoster}{HTML}{b0efff}
	\CabeceraDeSubSeccionConAnchoN{9}{0}{14}{80}{APLICACIÓN}{CTDiario}{CCPoster}
	\begin{NuevoParrafo}{11}{0}[16][16]
		\begin{Marco}[\LineaSupC][\LineaInfC][\LineaIzqC][\LineaDerC][CBlanco]
			\subseccionC{MetaModelo}
			\Diagrama{imgs/05-aplicacion/Meta}		
		\end{Marco}
	\end{NuevoParrafo}	
	\begin{NuevoParrafo}{11}{17}[17][16]
		\begin{Marco}[\LineaSupC][\LineaInfC][\LineaIzqC][\LineaDerC][CBlanco]
			\subseccionC{Uso de Aplicación}
			\Diagrama{imgs/05-aplicacion/UsoAplicacion}	
		\end{Marco}
	\end{NuevoParrafo}
	
	\begin{NuevoParrafo}{11}{34}[22][16]
		\begin{Marco}[\LineaSupC][\LineaInfC][\LineaIzqC][\LineaDerC][CBlanco]			
			\subseccionC{Estructura de Aplicación}
			\Diagrama{imgs/05-aplicacion/EstrAplicacion}		
		\end{Marco}
	\end{NuevoParrafo}	
	
	\begin{NuevoParrafo}{11}{56}[22][16]
		\begin{Marco}[\LineaSupC][\LineaInfC][\LineaIzqC][\LineaDerC][CBlanco]
			\subseccionC{Comportamiento de Aplicación}
			\Diagrama{imgs/05-aplicacion/CmtoAplicacion}		
		\end{Marco}
	\end{NuevoParrafo}	
		
	
	\begin{NuevoParrafo}{11}{78}[16][16]
		\begin{Marco}[\LineaSupC][\LineaInfC][\LineaIzqC][\LineaDerC][CBlanco]
			\subseccionC{Cooperación de Aplicación}
			\Diagrama{imgs/05-aplicacion/CoopAplicacion}	
		\end{Marco}
	\end{NuevoParrafo}


	\definecolor{CCPoster}{HTML}{82d44f}	
	\CabeceraDeSubSeccionConAnchoN{28}{0}{14}{80}{TECNOLOGÍA}{CTDiario}{CCPoster}
	\begin{NuevoParrafo}{30}{0}[30][15]
		\begin{Marco}[\LineaSupC][\LineaInfC][\LineaIzqC][\LineaDerC][CBlanco]
			\subseccionC{MetaModelo}
			\Diagrama{imgs/06-tecnologia/Meta}
		\end{Marco}
	\end{NuevoParrafo}
	
	\begin{NuevoParrafo}{30}{31}[17][12]
		\begin{Marco}[\LineaSupC][\LineaInfC][\LineaIzqC][\LineaDerC][CBlanco]
			\subseccionC{Tecnología}
			\Diagrama{imgs/06-tecnologia/Tecnologia}
		\end{Marco}
	\end{NuevoParrafo}
	
	\begin{NuevoParrafo}{30}{49}[25][12]
		\begin{Marco}[\LineaSupC][\LineaInfC][\LineaIzqC][\LineaDerC][CBlanco]
			\subseccionC{Uso de Tecnología}
			\Diagrama{imgs/06-tecnologia/UsoTecnologia}
		\end{Marco}
	\end{NuevoParrafo}
	
	\begin{NuevoParrafo}{46}{0}[13][9]
		\begin{Marco}[\LineaSupC][\LineaInfC][\LineaIzqC][\LineaDerC][CBlanco]
			\subseccionC{Realización del Servicio}
			\Diagrama{imgs/06-tecnologia/RealServicio}
			
		\end{Marco}
	\end{NuevoParrafo}
	
	\begin{NuevoParrafo}{46}{14}[16][9]
		\begin{Marco}[\LineaSupC][\LineaInfC][\LineaIzqC][\LineaDerC][CBlanco]
			\subseccionC{Físico}
			\Diagrama{imgs/06-tecnologia/Fisico}
		\end{Marco}
	\end{NuevoParrafo}
	
	\begin{NuevoParrafo}{42}{31}[17][13]
		\begin{Marco}[\LineaSupC][\LineaInfC][\LineaIzqC][\LineaDerC][CBlanco]
			\subseccionC{Despliegue e Implementación}
			\Diagrama{imgs/06-tecnologia/Implementacion}
		\end{Marco}
	\end{NuevoParrafo}
	\begin{NuevoParrafo}{42}{49}[25][13]
		\begin{Marco}[\LineaSupC][\LineaInfC][\LineaIzqC][\LineaDerC][CBlanco]			
			\subseccionC{Estructura de Información}
			\Diagrama{imgs/06-tecnologia/EstrInformacion}
		\end{Marco}
	\end{NuevoParrafo}
	
	\begin{NuevoParrafo}{30}{75}[19][25]
		\begin{Marco}[\LineaSupC][\LineaInfC][\LineaIzqC][\LineaDerC][CBlanco]			
			\subseccionC{Capas}
			\Diagrama{imgs/06-tecnologia/Capas}
		\end{Marco}
	\end{NuevoParrafo}

	\definecolor{CCPoster}{HTML}{f4deeb}
	\CabeceraDeSubSeccionConAnchoN{55}{0}{14}{53}{PROYECTO}{CTDiario}{CCPoster}
	\begin{NuevoParrafo}{57}{0}[22][8]
		\begin{Marco}[\LineaSupC][\LineaInfC][\LineaIzqC][\LineaDerC][CBlanco]
			\subseccionC{MetaModelo}
			\Diagrama{imgs/07-proyecto/Meta}
		\end{Marco}
	\end{NuevoParrafo}
	
	\begin{NuevoParrafo}{57}{20}[14][7]
		\begin{Marco}[\LineaSupC][\LineaInfC][\LineaIzqC][\LineaDerC][CBlanco]
			\subseccionC{Proyecto}
			\Diagrama{imgs/07-proyecto/Proyecto}
		\end{Marco}
	\end{NuevoParrafo}
	\begin{NuevoParrafo}{57}{35}[11][7]
		\begin{Marco}[\LineaSupC][\LineaInfC][\LineaIzqC][\LineaDerC][CBlanco]
			\subseccionC{Migración}
			\Diagrama{imgs/07-proyecto/Migracion}
		\end{Marco}
	\end{NuevoParrafo}
	\begin{NuevoParrafo}{57}{47}[20][7]
		\begin{Marco}[\LineaSupC][\LineaInfC][\LineaIzqC][\LineaDerC][CBlanco]
			\subseccionC{Imp/Migración}
			\Diagrama{imgs/07-proyecto/ImplementacionMigracion}
		\end{Marco}
	\end{NuevoParrafo}

	\begin{NuevoParrafo}{55}{68}[26][9]
	\begin{Marco}[\LineaSupC][\LineaInfC][\LineaIzqC][\LineaDerC][CBlanco]
		{\centering\FuenteTitulos\Large\textbf{Referencias}\par}
		\vspace{4pt}

		% Bibliografia de los .bib: hoy esta vacia porque nada usa \cite,
		% pero en cuanto se cite algo aparece aqui. heading=none para que
		% no imprima su propio titulo encima del de arriba.
		\renewcommand{\bibfont}{\normalfont\small}
		\printbibliography[heading=none]

		% Fuentes citadas a mano (no estan en los .bib).
		\definecolor{CURL}{HTML}{0000ff}
		\small
		\begin{enumerate}[label={[\arabic*]}, leftmargin=*, topsep=0pt, itemsep=4pt, parsep=0pt]
			\item OG, Archimate,
			{\color{CURL}\underline{\url{www.opengroup.org/archimate-forum}}},
			2024, Online, accessed \today
			\item Coloso,
			{\color{CURL}\underline{\url{www.colosoft.com.co}}},
			Online, accessed \today,\\
			Autor: Ph.D, M.Sc, Ing. Sistemas Sandro Bolaños, 2025.
		\end{enumerate}
	\end{Marco}
\end{NuevoParrafo}

\end{NuevaPagina}

\end{document}
